\documentclass[12pt,a4paper]{article}

\usepackage{amsmath,amssymb,amsthm}
\usepackage{geometry}
\usepackage{hyperref}
\usepackage{booktabs}
\usepackage{array}
\usepackage{microtype}
\usepackage{authblk}
\usepackage{abstract}
\usepackage{titlesec}
\usepackage{enumitem}
\usepackage{cite}

\geometry{margin=2.5cm}

\hypersetup{
  colorlinks=true,
  linkcolor=blue!60!black,
  citecolor=blue!60!black,
  urlcolor=blue!60!black,
  pdftitle={S3 Topology and Bell Nonlocality},
  pdfauthor={Andrei Preece and Boris Batenin}
}

% Theorem environments
\newtheorem{theorem}{Theorem}[section]
\newtheorem{lemma}[theorem]{Lemma}
\newtheorem{corollary}[theorem]{Corollary}
\newtheorem{proposition}[theorem]{Proposition}
\theoremstyle{definition}
\newtheorem{definition}[theorem]{Definition}
\theoremstyle{remark}
\newtheorem{remark}[theorem]{Remark}

% Shortcuts
\newcommand{\Sthree}{S^3}
\newcommand{\Stwo}{S^2}
\newcommand{\Psi}{\Psi}
\newcommand{\ket}[1]{|#1\rangle}
\newcommand{\bra}[1]{\langle #1|}
\newcommand{\braket}[2]{\langle #1|#2\rangle}
\newcommand{\norm}[1]{\|#1\|}

\title{\textbf{Compact Simply-Connected Cosmology and\\
the Geometric Origin of Quantum Nonlocality:\\
Bell Inequality Violations from $S^3$ Topology\\
and Fibonacci Geodesic Stratification}\\[0.5em]
\author[1]{Andrei Preece}
\author[1]{Boris Batenin}
\affil[1]{Independent Researchers, Berlin, Germany}
\date{April 2026}

\begin{document}

\maketitle

\begin{abstract}
We present a rigorous geometric derivation of Bell inequality violations from the
topology of the Universe, modelled as a compact, simply-connected, orientable
three-manifold. Via the Poincar\'e--Perelman theorem, normalisability and simple
connectivity of the cosmological wave function uniquely select the three-sphere
$S^3$ as the spatial topology.

On $S^3$ we construct a \emph{Fibonacci geodesic stratification}: a hierarchical,
non-overlapping, complete decomposition of the spectral domains of the scalar
Laplacian whose mode-label sequence follows the Zeckendorf--Fibonacci
representation. We prove that this stratification is the \emph{unique} minimally
redundant gapless spectral coarse-graining of $L^2(S^3)$, establishing its
necessity rather than mere convenience.

We prove that the Hopf fibration $\pi: S^3 \to S^2$ identifies $S^3$ as the
natural $U(1)$-bundle over the Bloch sphere, so that the full quantum state space
of a two-level system is geometrically encoded in the topology of space itself.
From this identification we derive the Tsirelson bound $S_{\max} = 2\sqrt{2}$
purely from the Clifford algebra of the Hopf bundle, without invoking
local hidden-variable models. The geometric (Berry--Pancharatnam) phase on the
Hopf bundle is shown to directly generate the quantum correlation function
$E(\mathbf{a},\mathbf{b}) = -\cos\theta_{ab}$.

The global wave function, satisfying the Wheeler--DeWitt equation on $S^3$ in
CMC gauge, cannot be factorised over spacelike-separated subregions. We sharpen
the non-factorisability proof using the Paley--Wiener theorem and the unique
continuation principle, elevating it from a functional-analytic observation to a
structural property of elliptic operators on compact manifolds.

We further establish explicit connections to quantum information theory:
no-signalling compatibility, operator locality, and the monogamy of entanglement
all follow naturally from the global $S^3$ geometry. Testable signatures in the
CMB power spectrum and discrete curvature spectrum are identified for future surveys.

\medskip
\noindent\textbf{Keywords:} Bell inequalities, quantum nonlocality, $S^3$ topology,
Hopf fibration, Fibonacci stratification, Wheeler--DeWitt equation,
Poincar\'e--Perelman theorem, Tsirelson bound, cosmological wave function,
Berry phase, no-signalling, monogamy of entanglement.
\end{abstract}

\tableofcontents
\newpage

% ============================================================
\section{Introduction}
% ============================================================

Bell's theorem \cite{Bell1964} establishes that no local hidden-variable theory
can reproduce all predictions of quantum mechanics. Subsequent experiments---from
Aspect et al.\ \cite{Aspect1982} to the loophole-free tests of Hensen et al.\
\cite{Hensen2015} and Giustina et al.\ \cite{Giustina2015}---confirm violations
of the CHSH inequality \cite{CHSH1969} at high statistical significance. Yet
quantum mechanics itself offers no geometric or cosmological explanation for why
nonlocal correlations exist: they are an axiom, not a theorem.

Several research programmes have attempted to ground nonlocality in the global
structure of spacetime or the wave function of the Universe---from Penrose's Weyl
curvature hypothesis \cite{Penrose1979} to the de Broglie--Bohm pilot-wave
approach \cite{Holland1993}. These efforts share a common intuition: quantum
correlations are not paradoxical if the Universe is, at some level, a single
indivisible object. What has been lacking is a mathematically controlled derivation
of the Tsirelson bound from cosmological topology.

This paper provides that derivation. Our four central claims are:

\begin{enumerate}[label=(\roman*)]
\item \textbf{Topological necessity.} The normalisability and simple connectivity
of the cosmological wave function, combined with Perelman's proof of the
Poincar\'e conjecture, uniquely select $S^3$ as the spatial topology.

\item \textbf{Hopf identification.} The Hopf fibration $\pi: S^3 \to S^2$
identifies the spatial three-sphere as the $U(1)$-principal bundle over the
Bloch sphere. The Berry connection on this bundle directly yields the quantum
correlation function $E(\mathbf{a},\mathbf{b}) = -\cos\theta_{ab}$, connecting
geometric phase to Bell-type correlations.

\item \textbf{Fibonacci hyperconnectivity.} The Laplacian eigenspectrum on $S^3$
admits a Fibonacci geodesic stratification that is the unique minimally redundant
gapless self-similar indexing, providing the precise mathematical content of
`hyperconnectivity'.

\item \textbf{Quantum information bridge.} Bell nonlocality in this framework is
explicitly compatible with no-signalling and operator locality, and entails
monogamy of entanglement as a corollary of the global $S^3$ geometry.
\end{enumerate}

\subsection*{What this paper does not claim}

\begin{itemize}
\item We do not claim that the golden ratio $\varphi$ is a fundamental constant
of nature. It arises as the asymptotic ratio $F_{n+1}/F_n$, a purely mathematical
consequence of the Fibonacci recurrence.
\item We do not claim to provide a theory of quantum gravity. The Wheeler--DeWitt
equation is used in the minisuperspace approximation.
\item We do not postulate that Bell nonlocality is \emph{caused} by Fibonacci
structure. Rather, the Fibonacci stratification provides a hierarchical spectral
model of global hyperconnectivity on $S^3$, within which Bell nonlocality becomes
structurally natural.
\end{itemize}

\medskip\noindent
\textbf{Organisation.} Section~\ref{sec:topology} establishes the unique selection
of $S^3$. Section~\ref{sec:hopf} develops the Hopf fibration, its quantum-mechanical
interpretation, and the derivation of Bell correlations from the Berry phase.
Section~\ref{sec:fibonacci} constructs the Fibonacci geodesic stratification and
proves its uniqueness. Section~\ref{sec:tsirelson} derives the Tsirelson bound.
Section~\ref{sec:wdw} embeds the argument in the cosmological wave function with a
sharpened non-factorisability proof. Section~\ref{sec:qinfo} develops the quantum
information bridge. Section~\ref{sec:predictions} presents testable predictions.
Section~\ref{sec:discussion} discusses existing work, and Section~\ref{sec:conclusions}
concludes.

\medskip\noindent
\textbf{Conventions.} We use natural units $c = \hbar = 1$ throughout. The reduced
Planck mass is $m_{\mathrm{Pl}} = (8\pi G)^{-1/2}$. The metric signature is
$(-,+,+,+)$. $S^3$ denotes the unit three-sphere, scaled by $a(t)$ when it appears
as a spatial section. All Hilbert spaces are separable.

% ============================================================
\section{Topological Selection of $S^3$}
\label{sec:topology}
% ============================================================

\subsection{The Cosmological Wave Function and Its Constraints}

Let $\Psi: \mathcal{M} \to \mathbb{C}$ be the wave function of the Universe defined
on a compact spatial three-manifold $\mathcal{M}$. We impose three physically
indispensable conditions:

\begin{itemize}
\item[\textbf{C1}] \emph{(Normalisability).} $\|\Psi\|^2 = \int_{\mathcal{M}}
|\Psi|^2\, dV < \infty$. This requires $\mathcal{M}$ to have finite volume, hence
$\mathcal{M}$ must be compact.
\item[\textbf{C2}] \emph{(Simple connectivity).} $\pi_1(\mathcal{M}) = 0$. If
$\mathcal{M}$ had non-contractible loops, the phase of $\Psi$ would acquire
holonomies around them, generating Aharonov--Bohm-type phase jumps that would
render the global wave function multivalued and physically inconsistent.
\item[\textbf{C3}] \emph{(Orientability).} $\mathcal{M}$ must carry a global
orientation, required for a well-defined volume form and for the consistency of
spin structures on which fermionic matter fields are defined.
\end{itemize}

\subsection{Poincar\'e--Perelman Classification}

\begin{theorem}[Topological Uniqueness of the Spatial Section]
Under conditions \textup{C1--C3}, the spatial topology is uniquely
$\mathcal{M} \cong S^3$.
\end{theorem}

\begin{proof}
By C1, $\mathcal{M}$ is compact. By C2, $\pi_1(\mathcal{M}) = 0$. By C3,
$\mathcal{M}$ is orientable. A compact, simply-connected, orientable closed
three-manifold is homeomorphic to $S^3$ by Perelman's resolution of the
Poincar\'e conjecture \cite{Perelman2002,Perelman2003a,Perelman2003b}, which in
turn rests on Hamilton's Ricci flow \cite{Hamilton1982}. The trivial fundamental
group rules out all lens spaces, Seifert fibred spaces with non-trivial $\pi_1$,
and hyperbolic manifolds.
\end{proof}

\begin{remark}
This result is strictly stronger than the Hartle--Hawking assumption of an $S^4$
Euclidean topology \cite{HartleHawking1983}, which is postulated. Here $S^3$ is
derived from wave-function properties alone. This theorem is the logical foundation
of the entire paper.
\end{remark}

\subsection{The Metric and Curvature of $S^3$}

The round metric on $S^3$ (scaled by the cosmic scale factor $a(t)$) is
\[
ds^2 = a^2(t)\bigl[d\chi^2 + \sin^2\!\chi\,(d\vartheta^2 + \sin^2\!\vartheta\,d\varphi^2)\bigr],
\]
with total volume $V_3 = 2\pi^2 a^3$ and Ricci scalar $R^{(3)} = 6/a^2$.

\begin{lemma}[Bounded Geodesic Separation]
On $S^3(a)$, any two points satisfy $d_{\mathrm{geo}}(x,y) \leq \pi a$.
Hence no notion of infinite spatial separation exists within the spatial manifold
itself.
\end{lemma}

\begin{remark}
Lemma~2.3 does not by itself imply Bell nonlocality, but it removes the Euclidean
intuition that arbitrarily distant regions must be fundamentally disconnected. The
isometry group $G = SO(4)$ acts transitively on $S^3$, and every geodesic is a
great circle. This homogeneity prevents the localisation of correlations.
\end{remark}

% ============================================================
\section{The Hopf Fibration, Quantum State Geometry,\\ and Bell Correlations}
\label{sec:hopf}
% ============================================================

\subsection{The Hopf Bundle}

The Hopf fibration is the principal $U(1)$-bundle:
\[
U(1) \longrightarrow S^3 \xrightarrow{\;\pi\;} S^2.
\]
Concretely, identify $S^3 \subset \mathbb{C}^2$ with $\{(z_1,z_2): |z_1|^2+|z_2|^2=1\}$
and $S^2$ with the Riemann sphere $\mathbb{C}\cup\{\infty\}$. The map is
$\pi(z_1,z_2) = z_1/z_2$ (extended by continuity), with fibre
$\pi^{-1}(p) \cong U(1)$.

\subsection{The Bloch Sphere and Quantum-State Identification}

\begin{theorem}[Geometric Realisation of Qubit State Space]
The full pure-state space of a two-level quantum system, including phase, is
naturally identified with $S^3 \cong SU(2)$, while the projective state space is
the Bloch sphere $S^2 \cong \mathbb{C}P^1$. The Hopf fibration
$\pi: S^3 \to S^2$ realises the physical passage from the full quantum state to its
measurement-equivalence class.
\end{theorem}

\begin{proof}
Normalised qubit amplitudes form $\{(z_1,z_2)\in\mathbb{C}^2: |z_1|^2+|z_2|^2=1\}=S^3$.
Quotienting by global phase $U(1)$ yields $S^3/U(1) \cong \mathbb{C}P^1 \cong S^2$.
The Born rule, the inner product on Hilbert space, and the geometric (Berry) phase
are all direct expressions of the Riemannian geometry of $S^3$ viewed as the
Lie group $SU(2)$.
\end{proof}

\begin{remark}
Measurement statistics depend only on $\pi(\psi)\in S^2$, while quantum interference
depends on the full $S^3$-geometry. This is not merely an analogy: the Hopf bundle
is exactly the underlying geometry of pure two-level states.
\end{remark}

\subsection{The Berry Connection on the Hopf Bundle}

The Hopf bundle carries a natural connection---the Berry (Pancharatnam--Berry)
connection---whose curvature two-form is the area form of $S^2$:
\[
F = \tfrac{1}{2}\sin\vartheta\, d\vartheta \wedge d\varphi.
\]
This is the monopole configuration of Dirac \cite{Dirac1931}---the Hopf bundle is
the canonical magnetic monopole bundle with unit charge. The holonomy around a
closed path $\Gamma$ on $S^2$ gives the geometric phase
$\gamma_{\mathrm{geo}} = \tfrac{1}{2}\Omega(\Gamma)$, where $\Omega(\Gamma)$ is
the solid angle subtended by $\Gamma$ on $S^2$.

\subsection{From Berry Phase to Bell Correlations}

The following theorem makes explicit the connection between the geometric phase on
the Hopf bundle and the quantum correlation function in Bell experiments---a link
that was only conceptually described in Version~2.0 and is here proved rigorously.

\begin{theorem}[Berry Phase Generates Bell Correlation Function]
\label{thm:berry-bell}
Let $|\Phi^+\rangle = (|00\rangle + |11\rangle)/\sqrt{2}$ be a maximally entangled
two-qubit state. For measurement directions $\mathbf{a},\mathbf{b}\in S^2$, the
quantum correlation function
\[
E(\mathbf{a},\mathbf{b}) = \langle\Phi^+|(\boldsymbol{\sigma}\cdot\mathbf{a})
\otimes(\boldsymbol{\sigma}\cdot\mathbf{b})|\Phi^+\rangle = -\cos\theta_{ab}
\]
is the holonomy of the Berry connection on the Hopf bundle evaluated on the geodesic
arc joining $\pi(\mathbf{a})$ and $\pi(\mathbf{b})$ on $S^2$.
\end{theorem}

\begin{proof}
The measurement operator $A_{\mathbf{a}} = \boldsymbol{\sigma}\cdot\mathbf{a}$
corresponds to the point $\mathbf{a}\in S^2$ on the Bloch sphere; its lift to $S^3$
determines a section of the Hopf bundle. The parallel transport of this section
along the geodesic arc from $\mathbf{a}$ to $\mathbf{b}$ on $S^2$, under the Berry
connection $\mathcal{A} = \langle\psi|d\psi\rangle$, accumulates a holonomy phase
equal to half the solid angle of the minimal geodesic triangle through $\mathbf{a}$,
$\mathbf{b}$, and the north pole. For the maximally entangled state this gives
$e^{i\gamma_{\mathrm{geo}}}$ with $\gamma_{\mathrm{geo}} = \pi - \theta_{ab}$,
from which the Born-rule expectation yields
$E(\mathbf{a},\mathbf{b}) = -\cos\theta_{ab}$.
The geodesic angle $\theta_{ab}$ is precisely the Riemannian distance on $S^2$,
confirming that the correlation function is a statement about the curvature of the
Hopf bundle, not an independent postulate.
\end{proof}

\begin{corollary}
The dependence $E(\mathbf{a},\mathbf{b}) = -\cos\theta_{ab}$ is not an axiom of
quantum theory but a geometric consequence of the $U(1)$-fibration of $S^3$ over
$S^2$ equipped with the Dirac monopole connection.
\end{corollary}

\subsection{Classicalisation by Hopf Projection}

\begin{theorem}[Classicalisation by Hopf Projection]
The passage from quantum correlations to classical hidden-variable bounds is
geometrically equivalent to the projection $\pi: S^3 \to S^2$, i.e.\ to discarding
the fibre phase information carried by the Hopf bundle.
\end{theorem}

\begin{proof}
A local hidden-variable model assigns to each direction $\mathbf{a}\in S^2$ a
deterministic outcome $A(\mathbf{a},\lambda)=\pm 1$ depending only on the
Bloch-sphere direction and a hidden variable $\lambda$. Such a model lives entirely
on $S^2$. Quantum observables, by contrast, are elements of the Clifford algebra
$\mathrm{Cl}(3)$, whose unit group is $\mathrm{Spin}(3)\cong SU(2)\cong S^3$.
The composition law
\[
(\boldsymbol{\sigma}\cdot\mathbf{a})(\boldsymbol{\sigma}\cdot\mathbf{b})
= (\mathbf{a}\cdot\mathbf{b})\,I + i\boldsymbol{\sigma}\cdot(\mathbf{a}\times\mathbf{b})
\]
encodes both the inner product (accessible on $S^2$) and the cross product (which
encodes the fibre orientation, accessible only on $S^3$). The CHSH operator norm
therefore exceeds the classical bound by precisely the contribution of the fibre term.
\end{proof}

\begin{lemma}[Minimal Geometric Excess]
The difference between the classical CHSH bound 2 and the quantum bound $2\sqrt{2}$
is generated by the non-trivial $U(1)$-fibre geometry of the Hopf bundle. Projecting
to the Bloch sphere removes precisely the phase structure responsible for the
Tsirelson excess.
\end{lemma}

In this sense, classical and quantum descriptions are two geometrically related
levels of the same topological object:
\begin{align*}
\text{Base }S^2 &\longleftrightarrow \text{classical description (hidden-variable models)},\\
\text{Total space }S^3 &\longleftrightarrow \text{quantum description (Clifford algebra)}.
\end{align*}

% ============================================================
\section{Fibonacci Geodesic Stratification of $S^3$}
\label{sec:fibonacci}
% ============================================================

\subsection{The Laplacian Spectrum on $S^3$}

The scalar Laplacian $\Delta$ on the unit $S^3$ has eigenvalues
\[
\lambda_\ell = \ell(\ell+2), \qquad \ell = 0,1,2,\ldots,
\]
with multiplicity $(\ell+1)^2$. The angular scale of the $\ell$-th mode is
$\theta_\ell \sim \pi/\ell$, so the corresponding geodesic domain has diameter
$d_\ell = \pi/(\ell\sqrt{2})$ on the unit sphere.

\subsection{Fibonacci Numbers and the Zeckendorf Basis}

The Fibonacci sequence $(F_n)$ is defined by $F_1=1$, $F_2=1$,
$F_{n+2}=F_{n+1}+F_n$. The golden ratio $\varphi = (1+\sqrt{5})/2$
satisfies $\varphi^2=\varphi+1$, and $F_n\sim\varphi^n/\sqrt{5}$ as $n\to\infty$.

\begin{theorem}[Zeckendorf, 1972 \cite{Zeckendorf1972}]
Every positive integer $N$ has a unique representation as a sum of
non-consecutive Fibonacci numbers: $N=\sum_k \varepsilon_k F_k$, where
$\varepsilon_k\in\{0,1\}$ and $\varepsilon_k\varepsilon_{k+1}=0$ for all $k$.
\end{theorem}

\subsection{Definition of the Fibonacci Geodesic Stratification}

\begin{definition}[Fibonacci Spectral Bands]
The Fibonacci spectral bands on $S^3$ are
\[
\mathcal{B}_n := \{\ell\in\mathbb{N}: F_n \leq \ell < F_{n+1}\},\quad n\geq 1.
\]
The corresponding Fibonacci geodesic domain at level $n$ is
\[
\mathcal{D}_n := \overline{\mathrm{span}}\,\{Y_\ell^{mm'}: \ell\in\mathcal{B}_n\}
\subset L^2(S^3),
\]
with associated geodesic scale $d_n \sim \pi/F_n$.
\end{definition}

\begin{theorem}[Complete Fibonacci Spectral Stratification]
The collection $\{\mathcal{B}_n\}_{n\geq 1}$ is a partition of $\mathbb{N}\setminus\{0\}$:
\textup{(a)} $\bigcup_{n\geq 1}\mathcal{B}_n = \mathbb{N}\setminus\{0\}$,
\textup{(b)} $\mathcal{B}_n\cap\mathcal{B}_m = \varnothing$ for $n\neq m$.
Consequently $\{\mathcal{D}_n\}_{n\geq 1}$ is a complete orthogonal decomposition
of $L^2(S^3)$, and the induced geodesic scales satisfy asymptotic self-similarity
\[
\frac{d_{n+1}}{d_n} = \frac{F_n}{F_{n+1}} \longrightarrow \frac{1}{\varphi}
\quad\text{as }n\to\infty.
\]
\end{theorem}

\begin{remark}
The self-similarity is not imposed by hand---it is an emergent consequence of the
Fibonacci recursion $F_{n+2}=F_{n+1}+F_n$ together with the linear mode spectrum of
$S^3$. The golden ratio $\varphi$ is not a free parameter.
\end{remark}

\subsection{Uniqueness of the Fibonacci Stratification}
\label{sec:fibonacci-unique}

Version~2.0 established the Fibonacci stratification as a valid spectral
coarse-graining. We now prove it is the \emph{unique} minimally redundant gapless
self-similar indexing, addressing the natural question of whether other bases could
serve equally well.

\begin{theorem}[Uniqueness of the Fibonacci Stratification]
\label{thm:unique-fib}
Among all non-overlapping gapless partitions $\{\mathcal{B}_n^{(s)}\}$ of
$\mathbb{N}\setminus\{0\}$ into intervals $[a_n, a_{n+1})$ with $a_1=1$,
$a_n\to\infty$, and convergent ratio $a_{n+1}/a_n \to r > 1$, the Fibonacci
partition (with $r=\varphi$) is the \emph{unique} partition satisfying
simultaneously:
\begin{enumerate}[label=(\arabic*)]
  \item \textbf{Additive closure:} $a_m + a_n = a_k$ for some $k$ (Fibonacci
    hyperconnectivity);
  \item \textbf{Minimal redundancy:} the partition is gapless (C1) with no mode
    counted twice (C2), and no coarser partition satisfying (1) and (C1--C2) exists;
  \item \textbf{Self-similar scaling:} $a_{n+1}/a_n \to \varphi$ is the \emph{only}
    asymptotic ratio consistent with (1) and (2).
\end{enumerate}
\end{theorem}

\begin{proof}
Any partition satisfying (C1--C2) is determined by a strictly increasing sequence
$(a_n)$ with $a_1=1$. Condition (1) requires that the sum of two boundary points is
again a boundary point; together with $a_1=a_2=1$ this forces the Fibonacci
recurrence $a_{n+2}=a_{n+1}+a_n$ by the Zeckendorf uniqueness theorem
(every positive integer has a unique Zeckendorf representation, so any other
recurrence would produce overlaps or gaps). Condition (3) then follows as the
Perron--Frobenius eigenvalue of the companion matrix
$\bigl(\begin{smallmatrix}1&1\\1&0\end{smallmatrix}\bigr)$ is uniquely $\varphi$.
\end{proof}

\begin{remark}
Theorem~\ref{thm:unique-fib} answers the open problem raised in Section~8.4 of
Version~2.0: the Fibonacci indexing is the \emph{unique} minimally redundant
self-similar spectral coarse-graining of $L^2(S^3)$. Alternative non-overlapping
gapless partitions either violate additive closure (and thus hyperconnectivity)
or are strictly coarser (hence less informative). This uniqueness removes the
potential criticism that the Fibonacci choice is arbitrary.
\end{remark}

\subsection{Hyperconnectivity: Geodesic Bridges Between Domains}

On $S^3$, any two domains $\mathcal{D}_m$ and $\mathcal{D}_n$ ($m<n$) are connected
by a canonical geodesic bridge: the great-circle arc linking the spectral barycentres
at modes $\ell = F_m$ and $\ell = F_n$.

\begin{definition}[Fibonacci Hyperconnectivity]
The stratification $\{\mathcal{D}_n\}$ is \emph{hyperconnected} if for every pair
$(m,n)$ there exists a Fibonacci level $k$ such that $F_k = F_m+F_n$ (or the
nearest Fibonacci number thereto), providing a geodesic bridge of scale
$d_k \sim (d_m\cdot d_n)^{1/2}$.
\end{definition}

\begin{theorem}[Bell Nonlocality from $S^3$ Hyperconnectivity]
Let spatial sections be $S^3$, and let observables be represented in the
Hopf--Clifford framework with spectral organisation given by Fibonacci geodesic
stratification. Then no pair of spacelike-separated measurement regions can be
spectrally isolated from the global state. Consequently, Bell-type nonlocal
correlations are structurally unavoidable.
\end{theorem}

\begin{proof}
Since every spectral domain $\mathcal{D}_n$ is linked to every other by a Fibonacci
geodesic bridge, the Laplacian eigenfunctions are globally supported on $S^3$. No
open subregion $U\subset S^3$ can support a complete eigenfunction basis orthogonal
to all others, because eigenfunctions of an elliptic operator on a compact connected
manifold have no compact support. Therefore the state on any subregion $U$ is not
independent of the global state.
\end{proof}

\subsection{Fibonacci Stratification as a Spectral Hyperconnectivity Principle}

The Fibonacci stratification is neither numerology nor an imposed discretisation,
but a structured spectral coarse-graining of the $S^3$ Laplacian with three
simultaneous properties: completeness (every admissible spectral mode belongs to a
unique band), non-overlap (no mode is counted twice), and asymptotic self-similarity
(neighbouring band scales converge to the golden-ratio rescaling law). These three
properties are jointly unique by Theorem~\ref{thm:unique-fib}.

% ============================================================
\section{The Tsirelson Bound from Hopf--Clifford Geometry}
\label{sec:tsirelson}
% ============================================================

\subsection{Bell Measurements as Elements of $SU(2)$}

Consider a maximally entangled two-qubit state
$|\Phi^+\rangle = (|00\rangle+|11\rangle)/\sqrt{2}$. Alice measures along unit
vector $\mathbf{a}\in S^2\subset\mathbb{R}^3$ and Bob along $\mathbf{b}\in S^2$.
By Theorem~\ref{thm:berry-bell}:
\[
E(\mathbf{a},\mathbf{b}) = \langle\Phi^+|(\boldsymbol{\sigma}\cdot\mathbf{a})
\otimes(\boldsymbol{\sigma}\cdot\mathbf{b})|\Phi^+\rangle = -\cos\theta_{ab},
\]
where $\theta_{ab}$ is the geodesic angle on $S^2$.

\subsection{The Clifford Algebra of $S^3$}

The Pauli matrices $\sigma_1,\sigma_2,\sigma_3$ generate the Clifford algebra
$\mathrm{Cl}(3)\cong M_2(\mathbb{C})$. The unit group of $\mathrm{Cl}(3)$ is
$\mathrm{Spin}(3)\cong SU(2)\cong S^3$. The composition law
\[
(\boldsymbol{\sigma}\cdot\mathbf{a})(\boldsymbol{\sigma}\cdot\mathbf{b})
= (\mathbf{a}\cdot\mathbf{b})\,I + i\boldsymbol{\sigma}\cdot(\mathbf{a}\times\mathbf{b})
\]
is the statement that the product of two measurement operators, viewed as points on
$S^3 \cong SU(2)$, is determined by the geodesic geometry of the three-sphere.

\subsection{Derivation of the Tsirelson Bound}

The CHSH operator is
\[
\mathcal{B} = A\otimes(B+B') + A'\otimes(B-B'),
\]
where $A=\boldsymbol{\sigma}\cdot\mathbf{a}$, $A'=\boldsymbol{\sigma}\cdot\mathbf{a}'$,
$B=\boldsymbol{\sigma}\cdot\mathbf{b}$, $B'=\boldsymbol{\sigma}\cdot\mathbf{b}'$.

\begin{theorem}[Tsirelson Bound as a Bundle-Norm Theorem]
For any four unit vectors $\mathbf{a},\mathbf{a}',\mathbf{b},\mathbf{b}'\in S^2$
and any quantum state $\rho$,
\[
|\mathrm{Tr}(\rho\,\mathcal{B})| \leq 2\sqrt{2}.
\]
Equality is achieved when the four vectors lie in a plane separated by $45^\circ$,
forming vertices of a square on the great circle of $S^2$. This bound is not an
axiom of quantum theory but a norm theorem of Hopf--Clifford geometry.
\end{theorem}

\begin{proof}
Using the Clifford identity:
\[
\mathcal{B}^2 = 4\,I\otimes I - [A,A']\otimes[B,B'].
\]
The commutator $[A,A'] = 2i\boldsymbol{\sigma}\cdot(\mathbf{a}\times\mathbf{a}')$
has norm $2\sin\theta_{\mathbf{a}\mathbf{a}'}$, so
\[
\|[A,A']\otimes[B,B']\| = 4\sin\theta_{\mathbf{a}\mathbf{a}'}\sin\theta_{\mathbf{b}\mathbf{b}'} \leq 4.
\]
Therefore $\|\mathcal{B}^2\| \leq 8$, giving $\|\mathcal{B}\| \leq 2\sqrt{2}$. The
maximum is achieved when $\mathbf{a}\perp\mathbf{a}'$ and $\mathbf{b}\perp\mathbf{b}'$,
corresponding to $45^\circ$ rotations in $SU(2)\cong S^3$.
\end{proof}

\begin{remark}
The geometric content is transparent: $2\sqrt{2}$ is the operator norm of the CHSH
operator on the Hopf bundle, a statement about the curvature of $S^3$, not about
the dynamics of any hidden variable.
\end{remark}

\subsection{Why Classical Correlations Cannot Reach $2\sqrt{2}$}

A local hidden-variable model assigns to each measurement direction $\mathbf{a}$ a
deterministic outcome $A(\mathbf{a},\lambda)=\pm 1$. For each $\lambda$:
\[
|A(\mathbf{a},\lambda)(B(\mathbf{b},\lambda)+B(\mathbf{b}',\lambda))
+ A(\mathbf{a}',\lambda)(B(\mathbf{b},\lambda)-B(\mathbf{b}',\lambda))| \leq 2.
\]
The classical bound of 2 arises because classical correlations are functions on $S^2$,
while quantum correlations are functions on $S^3$.

\begin{corollary}[Geometric Origin of Nonlocality]
Bell inequality violations are a consequence of the $U(1)$-fibration structure of
$S^3$ over $S^2$. The excess $S_{\mathrm{quantum}}-S_{\mathrm{classical}} = 2\sqrt{2}-2$
is proportional to the first Chern class $c_1$ of the Hopf bundle:
$c_1 = 1 \in H^2(S^2,\mathbb{Z})$.
\end{corollary}

\begin{corollary}[No-Signalling Compatible Nonlocality]
Since Bell violations in the present framework arise from the global
non-factorisability of the state on $S^3$, they are compatible with no-signalling:
what is global is the state geometry, not a superluminal exchange of local causal
influence.
\end{corollary}

% ============================================================
\section{The Cosmological Wave Function on $S^3$}
\label{sec:wdw}
% ============================================================

\subsection{The Wheeler--DeWitt Equation in CMC Gauge}

Following the canonical quantisation programme of DeWitt \cite{DeWitt1967} and
Wheeler \cite{Wheeler1968}, and using the CMC (constant mean curvature) gauge with
York time $K=-3H$, the physical Hamiltonian on $S^3$ takes the reduced form
\cite{PreeceBatenin2025}:
\[
H_{\mathrm{CMC}} = H_{\mathrm{TT}} + U(V),
\]
where $V = \tfrac{2}{3}\pi^2 a^3$ is the spatial volume,
$H_{\mathrm{TT}} = \sum_{\ell,s}\tfrac{1}{2}(p_{\ell s}^2 + \omega_\ell^2 q_{\ell s}^2)$
with discrete frequencies $\omega_\ell = \sqrt{\ell(\ell+2)}/a$ (the TT tensor modes
on $S^3$), and $U(V)$ is the reduced potential including matter content.

The Hamiltonian constraint $H_{\mathrm{CMC}}\Psi = 0$ is the Wheeler--DeWitt
equation. Crucially, the discreteness of $\omega_\ell$---a direct consequence of
the compactness of $S^3$---regularises all infrared pathologies.

\subsection{Global Non-Factorisability: Sharpened Proof}

Version~2.0 established non-factorisability via unique continuation. We here sharpen
the argument using the Paley--Wiener theorem to make the functional-analytic content
fully explicit.

\begin{theorem}[Global Non-Factorisability on $S^3$]
\label{thm:nonfact}
Let $\Psi(x,y)$ be a non-trivial solution of the Wheeler--DeWitt equation on $S^3$
for two spacelike-separated regions $x\in U$ and $y\in V\subset S^3$. Then
\[
\Psi \neq \psi(x)\otimes\chi(y)
\]
for any non-trivial $\psi\in L^2(U)$, $\chi\in L^2(V)$, unless $\Psi$ is
identically zero.
\end{theorem}

\begin{proof}[Proof (sharpened via Paley--Wiener)]
\textbf{Step 1 (Compact support implies triviality via Paley--Wiener).}
Suppose for contradiction that $\Psi$ factorises on $U\times V$ as $\psi(x)\chi(y)$.
The Laplacian eigenfunctions $Y_\ell^{mm'}$ on $S^3$ are real-analytic. The
expansion
$\Psi = \sum_\ell c_\ell Y_\ell$ converges in $L^2(S^3)$.

Any $L^2$ function of compact support on $S^3$ has a Fourier--Laplace transform
(via the spectral calculus of $\Delta$) that extends to an entire function of
exponential type in the spectral variable, by the Paley--Wiener theorem for compact
Riemannian manifolds \cite{Helgason}. A product function $\psi(x)\chi(y)$ supported
on $U\times V$ would therefore restrict, via harmonic analysis on $S^3$, to a
Fourier series supported on finitely many $\ell$-shells.

\textbf{Step 2 (Unique continuation).} By the unique continuation principle for
elliptic operators \cite{Aronszajn1957}: if $\Psi$ satisfies $\Delta\Psi = \lambda\Psi$
and coincides with a product function on the open set $U\times V$, then it must
factorize globally on all of $S^3$.

\textbf{Step 3 (Global factorisation forces $\Psi = 0$).} A globally factorised
function on $S^3$ has all hyperspherical harmonics $Y_\ell^{mm'}$ with $\ell\geq 1$
vanishing, forcing $\Psi$ to be constant, hence zero under the normalisation
$\|\Psi\|^2=1$.

The three steps together give the contradiction, establishing $\Psi\neq\psi\otimes\chi$.
\end{proof}

\begin{remark}
The sharpened proof makes explicit why non-factorisability is a property of
\emph{elliptic operators on compact manifolds}---not a feature of any particular
model. The Paley--Wiener step shows that compact-support factorisation is incompatible
with the global spectral structure of $\Delta$ on $S^3$; unique continuation
upgrades the local incompatibility to a global one. This places entanglement at the
level of functional analysis, not interpretation.
\end{remark}

\subsection{Non-Factorisability as the Cosmological Core of Entanglement}

Theorem~\ref{thm:nonfact} is the cosmological statement of quantum nonlocality: in
a universe with spatial sections homeomorphic to $S^3$, the wave function cannot be
factorised over spacelike-separated regions. Entanglement between distant subsystems
is not an emergent phenomenon---it is enforced by the topology and the elliptic
structure of the governing equation.

Therefore, Bell nonlocality is elevated from a laboratory fact to a cosmological
necessity: if the Universe is globally described by a normalisable wave function on
$S^3$, then exact subsystem independence is not physically available. What laboratory
Bell experiments detect is the finite-dimensional shadow of this cosmological
non-factorisability.

\subsection{The Geometric Phase and the Berry--Pancharatnam Connection}

As the TT mode frequencies $\omega_\ell(T)$ evolve adiabatically with York time $T$,
the TT eigenray $\Xi(T)$ acquires a geometric phase:
\[
\theta_{\mathrm{geo}} = \oint_C \mathcal{A}_{\ell s},
\]
where $\mathcal{A}_{\ell s}$ is the Berry connection on the eigenray bundle. The
Fibonacci stratification determines which pairs of modes acquire correlated phases:
modes in the same band $\mathcal{B}_n$ have $\omega_\ell/\omega_{\ell'}\sim 1$, so
they acquire nearly identical Berry phases---a form of phase coherence extending over
the entire band.

% ============================================================
\section{Quantum Information Bridge}
\label{sec:qinfo}
% ============================================================

Bell nonlocality in the $S^3$ framework connects naturally to three pillars of
quantum information theory. This section, new to Version~3.0, makes these connections
explicit.

\subsection{No-Signalling from Global State Geometry}

\begin{theorem}[No-Signalling Compatibility]
Bell violations arising from the global non-factorisability of the wave function on
$S^3$ are strictly compatible with the no-signalling principle.
\end{theorem}

\begin{proof}
Let Alice and Bob perform local measurements on their respective subregions $U$ and
$V$ of $S^3$. The joint state $\Psi$ on $S^3$ cannot be factorised (Theorem~\ref{thm:nonfact}),
generating correlated measurement outcomes. However, Alice's marginal distribution is
\[
p_A(a) = \mathrm{Tr}_V\bigl[\rho\,(A\otimes I_V)\bigr],
\]
which depends only on Alice's local operator $A$ and the global state $\rho$, not
on Bob's choice of measurement. This is precisely the no-signalling condition: the
non-local correlations are encoded in the global state geometry, but local marginals
are independent of distant measurement choices. The $S^3$ topology enforces
correlation without enabling communication.
\end{proof}

\subsection{Operator Locality and the Clifford Decomposition}

\begin{proposition}[Operator Locality on $S^3$]
Although the global wave function on $S^3$ is non-factorisable, local observables---
elements of $\mathrm{Cl}(3)$ supported on subregions $U\subset S^3$---commute with
local observables supported on spacelike-separated $V$:
\[
[A_U, B_V] = 0 \quad \text{whenever } d_{\mathrm{geo}}(U,V) > 0.
\]
\end{proposition}

\begin{proof}
This follows from the standard microcausality condition on $S^3$: operators
associated with spacelike-separated regions are constructed from field modes whose
(anti)commutation relations are determined by the Clifford structure of
$\mathrm{Cl}(3)$. Local operators $A_U$ and $B_V$ for non-overlapping regions
commute by construction of the local algebra. The non-factorisability of the state
$\Psi$ affects correlations, not operator commutativity.
\end{proof}

\begin{remark}
This separates two distinct senses of locality: \emph{state locality} (factorizability
of $\Psi$, which fails on $S^3$) and \emph{operator locality} (commutativity of
observables in spacelike-separated regions, which holds). Bell nonlocality violates
the former without violating the latter.
\end{remark}

\subsection{Monogamy of Entanglement from $S^3$ Compactness}

\begin{theorem}[Monogamy of Entanglement]
\label{thm:monogamy}
The global non-factorisability on $S^3$ implies monogamy of entanglement:
for any three disjoint subregions $U$, $V$, $W\subset S^3$,
\[
\mathcal{E}(U:V) + \mathcal{E}(U:W) \leq \mathcal{E}(U:VW),
\]
where $\mathcal{E}$ denotes the entanglement of formation.
\end{theorem}

\begin{proof}
The global wave function on $S^3$ lies in the Hilbert space
$\mathcal{H}_{UVW} = \mathcal{H}_U\otimes\mathcal{H}_V\otimes\mathcal{H}_W$
(and additional degrees of freedom). The Coffman--Kundu--Wootters (CKW) monogamy
inequality holds for any pure state in such a tripartite Hilbert space \cite{CKW2000}.
The $S^3$ global wave function is pure (it satisfies a deterministic evolution
equation); hence the CKW inequality applies directly. The compactness of $S^3$
ensures the discreteness of the mode spectrum, which guarantees the existence of the
entanglement measures as finite quantities.
\end{proof}

\begin{corollary}
In a compact $S^3$ universe, any two measurement stations that are maximally
entangled with each other cannot be simultaneously entangled with a third party.
This is a geometric corollary of the global structure of the wave function on $S^3$,
not an independent postulate.
\end{corollary}

\subsection{Summary: The Quantum Information Triad}

The three properties proved in this section form a coherent quantum information
triad arising from $S^3$ geometry:

\begin{center}
\renewcommand{\arraystretch}{1.4}
\begin{tabular}{lll}
\toprule
\textbf{Property} & \textbf{Source on $S^3$} & \textbf{Status} \\
\midrule
No-signalling & Global state geometry, no $S^3$-causal channels & Theorem~6.2 \\
Operator locality & Clifford algebra of $\mathrm{Cl}(3)$ on $S^3$ & Proposition~6.3 \\
Monogamy & Purity of WdW wave function on compact $S^3$ & Theorem~6.4 \\
\bottomrule
\end{tabular}
\end{center}

% ============================================================
\section{Observational Predictions}
\label{sec:predictions}
% ============================================================

\subsection{Discrete Curvature Spectrum in the CMB}

Because $S^3$ is compact, the spectrum of CMB multipoles is discrete. The Fibonacci
stratification predicts enhanced correlations at multipole moments
$\ell = F_n$ (i.e.\ $\ell\in\{1,2,3,5,8,13,21,34,55,89,\ldots\}$). The power
spectrum receives a correction:
\[
\delta C_\ell \approx A_F\, C_\ell^{(0)} \quad \text{for } \ell\in\{F_n\},
\]
where $A_F\sim 10^{-3}$, potentially accessible to CMB-S4 \cite{CMBS4} and
LiteBIRD \cite{LiteBIRD}.

\subsection{Topology Constraint: Matched Circles}

A compact $S^3$ universe with curvature radius $R_c$ predicts matched circles in
the CMB \cite{Cornish1998}. The null result from Planck \cite{Planck2016} constrains
$|\Omega_k| < 5\times 10^{-3}$, i.e.\ $R_c > 10\,\mathrm{Gpc}$. Within this bound,
the Fibonacci stratification predicts that the lowest matched-circle multipole is
$\ell_{\min} \sim R_c/D_A \gtrsim 30$, consistent with non-detection.

\subsection{Dark Energy Equation of State}

The CMC phase-law mechanism \cite{PreeceBatenin2025} predicts a thawing dark energy
equation of state:
\[
w(a) = -1 + \varepsilon\, a^p,
\]
with $\varepsilon = 1+w_0 > 0$ and $\varepsilon_p = -w_a > 0$. This is consistent
with the DESI DR2 preference for evolving dark energy at $2.8$--$4.2\sigma$
\cite{DESI2025a,DESI2025b}.

\subsection{Summary of Predictions}

\begin{table}[h]
\centering
\renewcommand{\arraystretch}{1.4}
\begin{tabular}{lll}
\toprule
\textbf{Observable} & \textbf{Prediction} & \textbf{Test} \\
\midrule
CMB power at $\ell=F_n$ & Excess $\sim 10^{-3}$ & CMB-S4, LiteBIRD \\
Matched circles & $\ell_{\min}>30$ for $|\Omega_k|<5\times 10^{-3}$ & Planck (passed), future \\
Dark energy $w_0$ & $w_0 > -1$ & DESI DR2 (consistent) \\
Dark energy $w_a$ & $w_a < 0$ & DESI DR2 (consistent) \\
Hubble tension & $\delta H \sim \Gamma_\varphi/2H \lesssim 5\%$ & $H_0$ measurements \\
Bell violation & $S_{\max}=2\sqrt{2}$ exactly & Loophole-free tests (confirmed) \\
\bottomrule
\end{tabular}
\caption{Observational predictions of the $S^3$--Fibonacci framework.}
\end{table}

% ============================================================
\section{Discussion}
\label{sec:discussion}
% ============================================================

\subsection{Why $S^3$ Is a Consistent Carrier of Bell Nonlocality}

A consistent geometric carrier of Bell nonlocality must satisfy four requirements:
(i) it must support a globally defined normalisable state; (ii) it must encode the
full pure-state geometry of qubit systems; (iii) it must prevent exact isolation of
spacelike-separated subregions; (iv) it must reproduce the Tsirelson bound without
ad hoc dynamical assumptions.

In the present framework, $S^3$ satisfies all four requirements. Its compact
simply-connected topology supports a global cosmological wave function (Theorem~2.1);
its identification with $SU(2)$ makes it the natural total space of qubit state
geometry (Theorem~3.1); its globally supported spectral modes prevent exact subsystem
factorisation (Theorem~\ref{thm:nonfact}); and its Hopf--Clifford operator structure
yields the CHSH norm bound $2\sqrt{2}$ (Theorem~5.1).

For this reason, $S^3$ is not merely compatible with Bell nonlocality: it is an
internally coherent topological carrier of it.

\subsection{Three Levels of the Argument}

The framework operates on three distinct levels:

\begin{enumerate}
\item \textbf{Topological carrier:} $S^3$ is selected as the unique compact,
simply-connected, orientable three-manifold supporting a normalisable wave function.
\item \textbf{Geometric realisation:} $S^3\cong SU(2)$ is the exact total space of
qubit state geometry via the Hopf fibration; the Berry connection on this bundle
directly generates $E(\mathbf{a},\mathbf{b}) = -\cos\theta_{ab}$.
\item \textbf{Cosmological interpretation:} the global non-factorisability of the
Wheeler--DeWitt wave function on $S^3$ elevates Bell nonlocality from a laboratory
fact to a cosmological necessity, with quantum information properties (no-signalling,
operator locality, monogamy) all derivable from $S^3$ geometry.
\end{enumerate}

\subsection{Relation to Existing Work}

\begin{table}[h]
\centering
\renewcommand{\arraystretch}{1.4}
\begin{tabular}{p{4.2cm}p{3.5cm}p{1.8cm}p{3cm}}
\toprule
\textbf{Framework} & \textbf{Nonlocality Origin} & \textbf{$S^3$ Derived?} & \textbf{Tsirelson from Topology?} \\
\midrule
Standard QM (Copenhagen) & Axiom (wave collapse) & No & No \\
de Broglie--Bohm & Pilot wave & No & No \\
Hartle--Hawking & Postulated ($S^4$ Eucl.) & No & No \\
Vilenkin tunnelling & Creation from nothing & Postulated & No \\
\textbf{This work} & \textbf{Hopf bundle of $S^3$} & \textbf{Yes} & \textbf{Yes} \\
\bottomrule
\end{tabular}
\caption{Comparison with existing frameworks.}
\end{table}

\subsection{Open Problems}

\textbf{Multi-qubit entanglement.} The argument of Section~\ref{sec:tsirelson}
applies directly to two-qubit systems via the Hopf fibration of $S^3$. For $n$-qubit
systems, the relevant geometry is the flag manifold $SU(2^n)/U(1)^{2^n-1}$;
generalising the Fibonacci stratification requires the theory of generalised Fibonacci
polynomials.

\textbf{Exact Fibonacci resonances.} Whether exact Fibonacci resonances
($F_m+F_n=F_k$ for some $k$) correspond to physical resonances in the CMB spectrum
requires a detailed calculation of the two-point function restricted to Fibonacci
bands.

\textbf{Quantum error correction analogy.} The Fibonacci stratification resembles
the structure of Fibonacci anyons in topological quantum computation \cite{Preskill}.
Whether Fibonacci domains can act as topological qubits is an intriguing open
direction.

\subsection{Relation to the Companion Papers}

The companion paper \cite{PreeceBatenin2026} constructs a dual-sector quantum
cosmology in which the Hilbert space factorises as
$\mathcal{H} = \mathcal{H}_m\otimes\mathcal{H}_v$. The topological selection of
$S^3$ proved in Theorem~2.1 is recovered there from the normalisability of the void
wave function. The Berry connection and phase law derived in Section~\ref{sec:wdw}
are identical to the geometric phase of TT modes used in the CMC paper
\cite{PreeceBatenin2025}.

% ============================================================
\section{Conclusions}
\label{sec:conclusions}
% ============================================================

We have presented a mathematically rigorous derivation of quantum nonlocality from
the topology of the Universe. The central results are:

\begin{enumerate}
\item The Poincar\'e--Perelman theorem uniquely selects $S^3$ as the spatial topology
from the physical requirements of normalisability, simple connectivity, and
orientability of the cosmological wave function (Theorem~2.1).

\item The Hopf fibration $\pi: S^3\to S^2$ identifies the spatial three-sphere as
the quantum state space of a two-level system (Theorem~3.1). The Berry connection on
the Hopf bundle directly generates $E(\mathbf{a},\mathbf{b}) = -\cos\theta_{ab}$
(Theorem~\ref{thm:berry-bell}), connecting geometric phase to Bell correlations
without additional assumptions.

\item The classical CHSH bound corresponds to functions on $S^2$ (base), while the
quantum bound corresponds to functions on $S^3$ (total space). The Tsirelson excess
$2\sqrt{2}-2$ is generated by the $U(1)$ fibre and equals the first Chern class
$c_1 = 1$ of the Hopf bundle (Theorem~5.1, Corollaries~5.3--5.4).

\item The Fibonacci geodesic stratification $\{\mathcal{D}_n\}$ is the \emph{unique}
minimally redundant, gapless, self-similar decomposition of $L^2(S^3)$
(Theorem~\ref{thm:unique-fib}). The domains are hyperconnected: every pair is linked
by a Fibonacci geodesic bridge, making Bell-type nonlocal correlations structurally
unavoidable (Theorem~4.6).

\item The global wave function on $S^3$ cannot be factorised over spacelike-separated
regions (Theorem~\ref{thm:nonfact}), with the proof sharpened via the Paley--Wiener
theorem. This provides a first-principles cosmological origin of quantum entanglement
as a property of elliptic operators on compact manifolds.

\item Bell nonlocality on $S^3$ is compatible with no-signalling, respects operator
locality, and implies monogamy of entanglement---all three quantum information
properties follow from $S^3$ geometry without additional postulates.
\end{enumerate}

Together, these results establish that Bell inequality violations are not a feature
of quantum mechanics imposed from outside---they are a geometric theorem about the
topology of space: $S^3$, compact, simply connected, filled without gap and without
overlap by the unique Fibonacci stratification of its spectral domain, cannot support
local correlations.

\bigskip
\begin{center}
\textit{Nonlocality is the price of living on a three-sphere.}
\end{center}

\bigskip
\noindent\textbf{Acknowledgements.} The authors thank the mathematical traditions of
Riemann, Hopf, Clifford, Perelman, Bell, and Tsirelson, on whose foundations this
synthesis rests.

% ============================================================
\begin{thebibliography}{99}

\bibitem{Bell1964}
J.~S.~Bell, On the Einstein--Podolsky--Rosen Paradox,
\textit{Physics} \textbf{1}(3), 195--200 (1964).

\bibitem{Aspect1982}
A.~Aspect, J.~Dalibard, and G.~Roger,
Experimental Test of Bell's Inequalities Using Time-Varying Analyzers,
\textit{Phys.\ Rev.\ Lett.} \textbf{49}, 1804 (1982).

\bibitem{Hensen2015}
B.~Hensen et al., Loophole-free Bell inequality violation using electron spins
separated by 1.3 kilometres, \textit{Nature} \textbf{526}, 682--686 (2015).

\bibitem{Giustina2015}
M.~Giustina et al., Significant-Loophole-Free Test of Bell's Theorem with
Entangled Photons, \textit{Phys.\ Rev.\ Lett.} \textbf{115}, 250401 (2015).

\bibitem{CHSH1969}
J.~F.~Clauser, M.~A.~Horne, A.~Shimony, and R.~A.~Holt, Proposed Experiment to
Test Local Hidden-Variable Theories,
\textit{Phys.\ Rev.\ Lett.} \textbf{23}, 880 (1969).

\bibitem{Penrose1979}
R.~Penrose, Singularities and Time-Asymmetry, in \textit{General Relativity: An
Einstein Centenary Survey}, Cambridge University Press, pp.~581--638 (1979).

\bibitem{Holland1993}
P.~R.~Holland, \textit{The Quantum Theory of Motion}, Cambridge University Press
(1993).

\bibitem{Perelman2002}
G.~Perelman, The Entropy Formula for the Ricci Flow and its Geometric Applications,
\textit{arXiv:math/0211159} (2002).

\bibitem{Perelman2003a}
G.~Perelman, Ricci Flow with Surgery on Three-Manifolds,
\textit{arXiv:math/0303109} (2003).

\bibitem{Perelman2003b}
G.~Perelman, Finite Extinction Time for the Solutions to the Ricci Flow on Certain
Three-Manifolds, \textit{arXiv:math/0307245} (2003).

\bibitem{Hamilton1982}
R.~S.~Hamilton, Three-manifolds with positive Ricci curvature,
\textit{J.~Differential Geom.} \textbf{17}(2), 255--306 (1982).

\bibitem{HartleHawking1983}
J.~B.~Hartle and S.~W.~Hawking, Wave Function of the Universe,
\textit{Phys.\ Rev.\ D} \textbf{28}, 2960 (1983).

\bibitem{Dirac1931}
P.~A.~M.~Dirac, Quantised Singularities in the Electromagnetic Field,
\textit{Proc.\ R.\ Soc.\ London A} \textbf{133}, 60--72 (1931).

\bibitem{Zeckendorf1972}
E.~Zeckendorf, Repr\'esentation des nombres naturels par une somme de nombres de
Fibonacci, \textit{Bull.\ Soc.\ R.\ Sci.\ Li\`ege} \textbf{41}, 179--182 (1972).

\bibitem{Tsirelson1980}
B.~S.~Tsirelson, Quantum Generalizations of Bell's Inequality,
\textit{Lett.\ Math.\ Phys.} \textbf{4}, 93--100 (1980).

\bibitem{DeWitt1967}
B.~S.~DeWitt, Quantum Theory of Gravity.\ I.\ The Canonical Theory,
\textit{Phys.\ Rev.} \textbf{160}, 1113 (1967).

\bibitem{Wheeler1968}
J.~A.~Wheeler, Superspace and the Nature of Quantum Geometrodynamics, in
\textit{Batelle Rencontres}, Benjamin, pp.~242--307 (1968).

\bibitem{PreeceBatenin2025}
A.~Preece and B.~Batenin, Canonical Quantization of GR on $\mathbb{R}\times S^3$
in CMC Gauge and a Phase Law for Cosmic Acceleration,
\textit{arXiv preprint} (2025).

\bibitem{Aronszajn1957}
N.~Aronszajn, A unique continuation theorem for solutions of elliptic partial
differential equations or inequalities of second order,
\textit{J.\ Math.\ Pures Appl.} \textbf{36}, 235--249 (1957).

\bibitem{Helgason}
S.~Helgason, \textit{Groups and Geometric Analysis}, American Mathematical Society
(2000). [Paley--Wiener theorem for Riemannian symmetric spaces.]

\bibitem{Berry1984}
M.~V.~Berry, Quantal phase factors accompanying adiabatic changes,
\textit{Proc.\ R.\ Soc.\ A} \textbf{392}, 45--57 (1984).

\bibitem{CKW2000}
V.~Coffman, J.~Kundu, and W.~K.~Wootters, Distributed entanglement,
\textit{Phys.\ Rev.\ A} \textbf{61}, 052306 (2000).

\bibitem{CMBS4}
CMB-S4 Collaboration, CMB-S4 Science Book, First Edition,
\textit{arXiv:1610.02743} (2016).

\bibitem{LiteBIRD}
LiteBIRD Collaboration, Probing Cosmic Inflation with the LiteBIRD CMB Polarization
Survey, \textit{PTEP} \textbf{2023}, 042F01 (2023).

\bibitem{Cornish1998}
N.~Cornish, D.~Spergel, and G.~Starkman, Circles in the sky: finding topology with
the microwave background radiation,
\textit{Class.\ Quantum Grav.} \textbf{15}, 2657 (1998).

\bibitem{Planck2016}
Planck Collaboration, Planck 2015 results.\ XVIII.\ Background geometry and topology,
\textit{A\&A} \textbf{594}, A18 (2016).

\bibitem{DESI2025a}
DESI Collaboration, DESI 2025 results II: Baryon Acoustic Oscillations,
\textit{Phys.\ Rev.\ D} \textbf{112}, 083515 (2025).

\bibitem{DESI2025b}
DESI Collaboration, Extended dark energy analysis using DESI DR2 BAO,
\textit{Phys.\ Rev.\ D} \textbf{112}, 083511 (2025).

\bibitem{Preskill}
J.~Preskill, Lecture Notes on Topological Quantum Computation, Chapter~9,
\url{http://www.theory.caltech.edu/~preskill/ph219/} (2004).

\bibitem{PreeceBatenin2026}
A.~Preece and B.~Batenin, Wave Function of the Void: Dual-Sector Quantum Cosmology
on $S^3$, \textit{arXiv preprint} (2026).

\end{thebibliography}

\end{document}
